\documentclass[11pt]{article}
\usepackage[margin=1in]{geometry}
\usepackage{amsmath,amssymb,amsthm}
\usepackage{graphicx}
\usepackage[hidelinks]{hyperref}
\setlength{\emergencystretch}{2em}

\newtheorem{theorem}{Theorem}
\newtheorem{lemma}{Lemma}
\newtheorem{corollary}{Corollary}
\newtheorem{remark}{Remark}

\newcommand{\C}{\mathbb C}
\newcommand{\dist}{\operatorname{dist}}

\title{BCP Forces a Faithful State and an Essential Blocker Ideal}
\author{Run \texttt{fa\_banach\_001}, lane 19}
\date{June 17, 2026}

\begin{document}
\maketitle

\section*{Source Problem}

The source paper is Jinghao Huang, Karimbergen Kudaybergenov, and Rui Liu,
\emph{The ball-covering property of von Neumann algebras and noncommutative
symmetric spaces}, arXiv:2311.14257.  After proving the commutative unital
criterion for \(C(K)\), it asks for a noncommutative \(C^*\)-algebra
characterization:

\begin{center}
\includegraphics[width=0.92\textwidth]{figures/open_problem_crop.png}
\end{center}

\section*{Background}

We use two previous lane-19 packets.

First,
\[
\texttt{2311.14257\_bcp\_countable\_capture\_exact\_equivalence}
\]
proved that a normed space fails BCP if and only if it has countable capture.
Equivalently, a normed space \(X\) has BCP if and only if there is a
countable set \(C\subset X\) such that for every \(u\in S_X\) some
\(c\in C\) satisfies
\[
        \dist(c,\C u)<\|c\|.
\]

Second,
\[
\texttt{2311.14257\_separable\_essential\_ideal\_ubcp\_reduction}
\]
proved that every \(C^*\)-algebra containing a separable essential ideal has
UBCP.  The goal of this packet is to test how close BCP itself comes to
producing such an ideal.

\section*{Faithful State Consequence}

\begin{theorem}[BCP gives a faithful state]
\label{thm:faithful}
Let \(A\) be a \(C^*\)-algebra.  If \(A\) has BCP, then \(A\) admits a
faithful state.
\end{theorem}

\begin{proof}
The zero algebra is trivial, so assume \(A\ne0\).  Let
\[
        S_A\subseteq\bigcup_{n=1}^{\infty}B(c_n,r_n)
\]
be a BCP cover by open balls, with \(r_n<\|c_n\|\).  Ignore any redundant
zero centres.  For each \(n\), choose a state \(\varphi_n\) on the unitization
of \(A\) such that
\[
        \varphi_n(c_n^*c_n)=\|c_n\|^2.
\]
This is possible because states norm positive elements.

Define bounded linear functionals on \(A\) by
\[
        f_n(x)=\varphi_n(c_n^*x).
\]
We claim that the family \((f_n)\) separates points of \(A\).  If not, choose
nonzero \(x\in A\) with \(f_n(x)=0\) for every \(n\), and put
\(u=x/\|x\|\).  Let
\[
        (\pi_n,H_n,\xi_n)
\]
be the GNS representation of \(\varphi_n\).  Then
\[
        \langle \pi_n(u)\xi_n,\pi_n(c_n)\xi_n\rangle
        =
        \varphi_n(c_n^*u)
        =
        0.
\]
Hence
\[
\begin{aligned}
        \|u-c_n\|^2
        &\ge
        \|\pi_n(u-c_n)\xi_n\|^2                                      \\
        &=
        \|\pi_n(u)\xi_n\|^2+\|\pi_n(c_n)\xi_n\|^2
        \ge
        \|c_n\|^2.
\end{aligned}
\]
Therefore \(\|u-c_n\|\ge\|c_n\|>r_n\) for every \(n\), contradicting that the
balls cover \(S_A\).

Thus the \(f_n\)'s separate points.  Put
\[
        \psi=\sum_{n=1}^{\infty}2^{-n}\varphi_n|_A .
\]
After normalizing if necessary in the nonunital case, \(\psi\) is a state on
\(A\).  We show it is faithful.  If \(a\in A\) and \(\psi(a^*a)=0\), then
\[
        \varphi_n(a^*a)=0
        \qquad(n\ge1).
\]
In the GNS representation, this means \(\pi_n(a)\xi_n=0\).  Hence
\[
        f_n(a)=
        \langle \pi_n(a)\xi_n,\pi_n(c_n)\xi_n\rangle
        =
        0
        \qquad(n\ge1).
\]
Since the \(f_n\)'s separate points, \(a=0\).  Thus \(\psi\) is faithful.
\end{proof}

\begin{corollary}
Every \(C^*\)-algebra without a faithful state fails BCP.
\end{corollary}

\section*{Essential Ideal Generated by a Blocker}

\begin{theorem}[BCP blockers generate essential ideals]
\label{thm:essential-generated}
Let \(A\) be a nonzero \(C^*\)-algebra with BCP, and let
\(C\subset A\) be a countable set such that for every \(u\in S_A\) there is
\(c\in C\) with
\[
        \dist(c,\C u)<\|c\|.
\]
Let \(I\triangleleft A\) be the closed two-sided ideal generated by \(C\).
Then \(I\) is essential.
\end{theorem}

\begin{proof}
Suppose \(I\) is not essential.  Then there is a nonzero closed two-sided
ideal \(J\triangleleft A\) with \(I\cap J=\{0\}\).  For ideals in a
\(C^*\)-algebra, this implies
\[
        IJ=JI=\{0\}.
\]
Choose \(u\in J\) with \(\|u\|=1\).

For every \(c\in C\) and every scalar \(\lambda\), the element \(c-\lambda u\)
lies in the orthogonal \(C^*\)-direct sum \(I\oplus J\).  Therefore
\[
        \|c-\lambda u\|
        =
        \max\{\|c\|,|\lambda|\,\|u\|\}
        \ge
        \|c\|.
\]
Taking the infimum over \(\lambda\in\C\) gives
\[
        \dist(c,\C u)=\|c\|
        \qquad(c\in C).
\]
This contradicts the assumed blocker property of \(C\).  Hence \(I\) is
essential.
\end{proof}

\section*{Relation to the UBCP Bridge}

Combining Theorem~\ref{thm:essential-generated} with the exact
BCP/countable-capture equivalence gives:

\begin{corollary}
\label{cor:bridge}
Every \(C^*\)-algebra with BCP has a countable subset whose generated ideal is
essential.
\end{corollary}

This is close to the previous UBCP theorem, but it is not the same as having
a separable essential ideal.  The distinction is real:

\begin{remark}[why countable generation is not enough]
\label{rem:rankone}
Let \(H\) be a nonseparable Hilbert space and let \(p\) be a rank-one
projection in \(\mathcal B(H)\).  The closed ideal of \(\mathcal B(H)\)
generated by \(p\) is \(\mathcal K(H)\), which is essential in
\(\mathcal B(H)\) but nonseparable.  Thus an ideal generated by one element
inside a larger \(C^*\)-algebra need not be separable.

This example does not give BCP: \(\mathcal B(H)\) fails BCP by the von
Neumann algebra classification in the source paper.  It only shows that the
last missing bridge cannot be reduced to the formal phrase ``countably
generated ideal.''
\end{remark}

\section*{What Remains}

The packet
\[
\texttt{2311.14257\_separable\_essential\_ideal\_ubcp\_reduction}
\]
proves:
\[
        \text{separable essential ideal}
        \quad\Longrightarrow\quad
        \text{UBCP}.
\]
The present packet proves:
\[
        \text{BCP}
        \quad\Longrightarrow\quad
        \text{a countable blocker generates an essential ideal}.
\]
The still-open bridge is whether the special geometry of a BCP blocker forces
enough separable local structure inside that essential ideal to construct a
uniform countable cover.  Equivalently, a hypothetical BCP-not-UBCP
\(C^*\)-algebra would have a faithful state and an essential ideal generated
by a countable BCP blocker, but no separable essential ideal of the kind used
in the previous positive theorem.

\section*{Novelty and Search Bounds}

Local run indexes were searched for arXiv:2311.14257, BCP, UBCP, faithful
state, weak-star separability, essential ideal, countable blocker, and
separable essential ideal.  Web searches on June 17, 2026 for
``ball-covering property faithful state C*-algebra'' and variants did not
locate this exact C*-specific consequence.  The density-continuum packet
already recorded a Banach-space countable-separating-functional consequence
of BCP; Theorem~\ref{thm:faithful} upgrades that mechanism to a faithful state
in the C*-algebra setting.

\section*{Human Review Focus}

Reviewers should check:
\begin{enumerate}
    \item the GNS Pythagorean estimate in Theorem~\ref{thm:faithful};
    \item that the resulting functionals really separate all points of \(A\);
    \item the faithfulness of the convex combination of the states
          \(\varphi_n\);
    \item the orthogonal direct-sum norm calculation in
          Theorem~\ref{thm:essential-generated};
    \item the limitation in Remark~\ref{rem:rankone}.
\end{enumerate}

\begin{thebibliography}{9}

\bibitem{HKL}
Jinghao Huang, Karimbergen Kudaybergenov, and Rui Liu,
\emph{The ball-covering property of von Neumann algebras and noncommutative
symmetric spaces},
arXiv:2311.14257, 2023.

\bibitem{LLLZ}
Minzeng Liu, Rui Liu, Jimeng Lu, and Bentuo Zheng,
\emph{Ball covering property from commutative function spaces to
non-commutative spaces of operators},
arXiv:2111.04921, 2021.

\bibitem{Blackadar}
Bruce Blackadar,
\emph{Operator Algebras: Theory of \(C^*\)-Algebras and Von Neumann
Algebras},
Encyclopaedia of Mathematical Sciences 122, Springer, 2006.

\end{thebibliography}

\end{document}
